% Folder 83, batch 08 — archivists' pages 141–160.
% Pass: Opus 5.5 (claude-opus-5-5), 2026-10-10 — first pass, unchecked against the pages by a human.
%
% Very fast hand throughout: the formulas come through, much of the connecting prose does not.
\documentclass[11pt,a4paper]{article}
\input{../preamble/grothendieck}

\folder{83}
\batch{8}
\pages{141}{160}
\foldertitle{[Icosaèdres, bi-icosaèdres et algèbres d'Azumaya de rang 4] : notes manuscrites (1982-1983, 1986, s.d.).}
\dating{1982-1986}
\watermark{Édition de démonstration}

\begin{document}

\page{141}
\note{la page reprend au milieu d'un argument commencé avant ce lot : $T$, $T'$, $\delta$, $L$, $L'$, $L_P$, $b$, $b'$ y sont déjà en place.}

\[
  \delta(T) = \delta(T')
\]
donc $T \mapsto T'$ induit un \struck{\ill{}} automorphisme de
\[
  \widetilde{S} = X \wedge X' = \delta^{-1}\bigl(\text{section nulle de } \mathbb{V}(L \otimes L')\bigr).
\]
Quand on choisit bases $u$, $v$ dans $L$\note{lecture du premier facteur incertaine ; peut-être $H$, $H'$.} \ill{}, $L_P$ identifié à l'axe des $t$, l'involution $T \mapsto T'$ \struck{équi} se représente par
\[
  t \longmapsto bb' - t ,
\]
qui \struck{$\Rightarrow$} \add{induit} \emph{l'involution can.} sur $\widetilde{S}$, schéma des solutions de $t^2 - bb't + {}$ \ill{} $= 0$.

\marginal{$\triangle$} On va donner une interprétation schématique de $X \wedge X'$ comme un produit contracté.

\textbf{Prop.} Supposons que $2$ soit inv., ou plus généralement, qu'en tout pt où $2$ est nul, soit $X$ soit $X'$ est étale (cela signifie aussi que $2$, $\delta$, $\delta'$ ne s'annulent pas simultanément sur $S$, i.e.\ que l'idéal $(2, \delta, \delta')$ dans $k$ est l'idéal unité).\note{sous $\delta$ et $\delta'$ il écrit $\delta = b^2 - 4c$, $\delta' = b'^2 - 4c'$.} Alors le \struck{\ill{}} morphisme
\[
  X \times_S X' \longrightarrow X \wedge_S X'
\]
est \struck{\uncertain{fid}} un morphisme quadratique (fini loc.\ libre de rang $2$) et, si $\sigma_X$, $\sigma_{X'}$ sont les inv.\ can.\ de $X$ et $X'$ resp., $\sigma_X \times_S \sigma_{X'}$ est l'inv.\ can.

\page{142}
de $X \times_S X'$ sur $X \wedge_S X'$, et celui-ci est le \struck{sous} schéma quotient (universel)
\[
  X \wedge_S X' \simeq X \times_S X' \big/ \underbrace{\sigma_X \times_S \sigma_{X'}}_{= \sigma_P}
\]
conformément au yoga général des produits contractés. Pour que $X \times_S X' \to X \wedge_S X'$ soit étale, il faut et il suffit qu'en tout $s \in S$, $X_s$ \emph{ou} $X'_s$ soit étale. Pour que $X \wedge_S X'$ soit étale sur $S$, il faut et il suffit que \struck{\ill{} tout point selon}, $X$ \add{et} $X'$ soient étales. (Dans ce cas, on retrouve la notion familière. Si $X$ seule est étale, on peut considérer $X \wedge X'$ comme un « twisté » de $X'$ par le $(\mathbb{Z}/2\mathbb{Z})_S$-torseur $X$ et l'op.\ de $\mathbb{Z}/2\mathbb{Z}$ sur $X'$ par son inv.\ can.)

En termes des alg.\ $H$, $H'$ de $X$, $X'$, et $P = H \otimes H'$ de $X \times_S X'$ [\uncertain{non cette fois} la structure d'alg.\ commut.\ can.\ sur $P$ !] on peut décrire celle de $X \wedge X'$ comme la sous-alg.\ \add{$H \wedge H'$} des invariants sous \ill{}. Si $u$, $v$ sont des bases, donc $H \wedge H'$ a comme base $1$ ($= 1 \otimes 1$) et
\[
  \boxed{\; u * v \overset{\mathrm{def}}{=} \operatorname{Tr}(v)\,u + \operatorname{Tr}(u)\,v - 2\underbrace{uv}_{= u \otimes v} \in (H \otimes H')^{\sigma_P} \;}
\]

\marginal{NB Cet élément \ill{} (défini \ill{} condition \ill{} sur $H$, $H'$ \ill{}) donc \ill{}}

\page{143}
\[
  \left\{
  \begin{aligned}
    \operatorname{Tr}(u * v) &= bb' \\
    \det(u * v) &= \delta(u, v) = b^2c + b'^2c' - 4cc'
  \end{aligned}
  \right.
\]
\note{les accents sur $b$, $c$ dans la seconde ligne sont mal visibles ; on transcrit ce qu'on lit.}

\marginal{il s'agit ici de la trace et du déterminant en tant qu'éléments d'une alg.\ quadratique [$H \wedge H'$], sur $S$}

La formule \struck{donnée} plus \struck{\ill{}} $T_{\varphi,\varphi'}(u,v)$ peut s'écrire
\[
  \boxed{\; T_{\varphi,\varphi'}(u,v) = \varphi \otimes \varphi'(u * v) \;}
\]

Notons aussi les formules
\[
  \underset{\text{discr}}{\delta(u * v)} = \delta(u)\,\delta(v)
\]
ainsi que
\[
  \left\{
  \begin{aligned}
    (u * v)' &= u' * v = u * v' \\
    u' * v' &= u * v
  \end{aligned}
  \right.
\]
\note{sous $(u*v)'$ : « le $'$ étant relatif à l'alg.\ quad.\ $H \wedge H'$ ».}

\textbf{NB} Toutes les formules précédentes ont un sens et sont valables, sans restriction de \uncertain{cas} sur $S$, $H$, $H'$, pour \emph{tt} $u \in H$ et $v \in H'$ (base ou pas !)

On récupère $L_P \overset{\mathrm{df}}{=} L_{H,H'}$ à partir de $\widetilde{S} = X \wedge X'$ comme l'espace affine enveloppant (au-dessus de $S$) du rev.\ quad.\ $\widetilde{S}$, et
\[
  \delta : L_P \longrightarrow (L \otimes L')^{-1}
\]
comme l'hom (discriminant !) canonique associé à \struck{\ill{}} ce revêtement.

\marginal{NB Si $u$, $v$ sont des bases, \ill{} $u * v$ \ill{} et \ill{} dans \ill{} autre que $\widetilde{S} \subset \mathbb{V}(L_P)$ \ill{} … Ainsi la fonction $t_{u,v}$ sur $\mathbb{V}(L_P)$ \ill{} considérée comme \ill{} $H^\vee \otimes H'^\vee$, \ill{} $\in H \otimes H'$ \ill{} restriction \ill{} à $L_P$}
\note{longue note en diagonale dans l'angle supérieur gauche, très serrée ; seuls quelques symboles se lisent.}

\page{144}
\note{les deux premières lignes achèvent la phrase qui termine le bas de la page 145 (« donc $k[v] = H'$ et $k[w] = H''$ engendrent sous-alg. ») ; les deux feuillets semblent intervertis dans la liasse.}
\[
  k + k\cdot v + kw
\]
qui est de rang $3$, et pas égal à $P$…

Je voudrais expliciter tout dans le cas où $H$ et $H'$ sont \uncertain{unipotentes}, plus précisément
\[
  H = k \oplus L, \qquad H' = k \oplus L', \qquad
  P = H \otimes H' = \underbrace{k \oplus L \oplus L'}_{} + L \otimes L'
\]
avec $L$, $L'$ formés d'éléments de carré nul.

\marginal{i.e.\ $X$, $X'$ sont définis comme schémas des \ill{} \struck{\ill{}} de $L \to L^{\otimes 2}$, $L' \to L'^{\otimes 2}$ triviales}

Donc \struck{\ill{}} $T_{H,H'}$ s'identifie \struck{à} à $(L \otimes L')^\vee$. \struck{Choisir \add{$L$} base $u \in L$, $v \in L'$},
\[
  \delta_{H,H'} : \underset{\displaystyle\simeq\ (L \otimes L')^{-1}}{T_{H,H'}} \longrightarrow (L \otimes L')^{\otimes -2}
\]
\[
  (L \otimes L')^{-1} \xrightarrow[\text{élévation au carré}]{} (L \otimes L')^{\otimes -2}
\]
donc
\[
  H \wedge H' \simeq k \oplus (L \otimes L') \qquad \text{($L \otimes L'$ formé d'élts de carré nul)}
\]

\struck{Si $u \in L$, $v \in L'$}

L'hom
\[
  X \times_S X' \longrightarrow X'' = X \wedge X'
\]
est donné par
\[
  \underset{\substack{\text{sections de } L, L' \\ \text{de carré nul}}}{(\xi, \eta)} \longrightarrow \underset{\substack{\text{section de } L \otimes L' \\ \text{de carré nul}}}{\xi \otimes \eta}
\]
L'hom d'alg.\ correspondant
\[
  H \otimes H' \simeq k \oplus L \oplus L' \oplus L \otimes L' \longleftarrow H'' = k \oplus L \otimes L'
\]
est celui défini par
\note{la phrase reste en suspens au bas de la page.}

\page{145}
\textbf{NB} On a, \struck{pour} la structure d'Alg.\ canonique \add{commut.} (produit tens.) sur $P = H \otimes H'$, trois sous-alg.\ quadratiques
\[
  H, \quad H', \quad H''
\]
($P/k \simeq H/k \oplus H'/k \oplus H''/k$ comme modules) et trois involutions mutuellement commutantes
\[
  \sigma_H, \quad \sigma_{H'} \quad \text{et} \quad \sigma_{H''}
\]
de telle façon que
\[
  H = P^{\sigma_H}, \quad H' = P^{\sigma_{H'}}, \quad H'' = P^{\sigma_{H''}}, \qquad
  H \cap H' = H' \cap H'' = H'' \cap H = k \quad (\text{\uncertain{union} !})
\]

\marginal{(où $H'' = H \wedge H'$ --- on suppose maintenant $(2, b, b')k = k$) \quad ($\sigma_{H''}$ défini comme $\sigma_H \sigma_{H'}$)}

\struck{et bien que $P$ soit \ill{}} $P$ est \add{engendré comme} alg.\ par $H$ et $H'$ (d'où $H \otimes H' \xrightarrow{\sim} P$), \struck{et $P$ est aussi} \ill{} $H \wedge H'$). (Ce qui caractérise $H''$ comme $H \wedge H'$.) Pour que $P$ soit aussi engendré par $H$ et $H''$ \ill{} par $H'$ et $H''$, il faut et il suffit que \ill{} $H$ \ill{}, $H'$ (et alors \ill{} $H''$ \ill{}). Si \ill{} alors $P$ est engendré par $H, H'$ et $H, H''$, \ill{} entre $H'$ et $H''$ \ill{} : $H$ \ill{} (La relation
\[
  H'' \simeq H \wedge H' \iff H' \simeq H \wedge H''
\]
symétrique en $H'$ et $H''$), mais \emph{non} pour $H'$ et $H''$. En effet : prenant une base $u$ de $H$, $v$ de $H'$ \ill{} $H'$ \ill{}
\[
  u^2 = u, \quad v^2 = 0 \qquad (\text{donc } b = 1,\ c = b' = c' = 0)
\]
et posant
\[
  w = u * v = v - 2uv
\]
on a \struck{bien} \add{base de $P$} \struck{$w^2 = 0$, $vw = $}
\[
  1, \quad u, \quad v, \quad w
\]
mais aussi
\[
  v^2 = w^2 = vw = 0
\]
donc $k[v] = H'$ et $k[w] = H''$ engendrent sous-alg.

\page{146}
\note{la page continue la phrase laissée en suspens au bas de la page 144.}
\[
  u \underset{L,L'}{\otimes} v \longmapsto u * v = -2uv
\]
C'est \uncertain{nul} ! \uncertain{réguliè} comme il se doit, quand $2$ est inv. Par contre, en car.\ $2$ \struck{c'est identique} (i.e.\ \ill{} sur \struck{$P$} $H$, $H'$, $H''$, sont identiques : l'identité !) c'est id \ill{}. D'autre part, $H''$ ne s'identifie pas, donc, à une sous-algèbre de $H \otimes H'$.

\marginal{Plus généralement, \ill{} des \ill{} $\mathcal{X}(L, q)$ \ill{} par des $L$ \ill{}, $\mathcal{X}(L, q) \wedge \mathcal{X}(L', q')$ \ill{} (iso.\ can.) \ill{} $\simeq \mathcal{X}(L \otimes L', q \otimes q')$ \ill{} \ill{} si $2$ \ill{} \ill{} des \ill{} \ill{} \ill{}…}
\note{note en diagonale dans la marge droite ; seuls les symboles se lisent avec quelque sûreté.}

Soit $X_L = \operatorname{Spec} k \oplus L$, \uncertain{\ill{}} \ill{} en carré nul. \marginal{Signe contr.\ \ill{} du car.\ \ill{}} On a donc
\[
  X_L \wedge X_{L'} \simeq X_{L \otimes L'} .
\]
Ceci montre que les propriétés du produit contracté sont assez différentes de celles des \struck{produit} pr.\ V.\ \ill{} familières pour les $(\mathbb{Z}/2\mathbb{Z})_S$-torseurs. Ainsi, le donnée
\[
  X_{L''} \simeq X_L \wedge X_{L'} \qquad \text{i.e.\ } L'' \simeq L \otimes L'
\]
n'est pas \ill{} \ill{} symétrique \ill{} $L, L', L''$), et
\[
  X_L \wedge X_L \simeq \struck{\ill{}} X_{L^{\otimes 2}} \simeq k + L^{\otimes 2}
\]
n'est pas \uncertain{com.}\ \uncertain{neutre} : $(\mathbb{Z}/2\mathbb{Z})_S$ ! [Par contre, si $X$ quadratique sur $S$, $X \wedge X$ \ill{} si $X$ \ill{}, et \ill{} bien \ill{} \ill{} $X \wedge X \simeq (\mathbb{Z}/2\mathbb{Z})_S$ …].

\page{147}
Il faudrait aussi examiner l'associativité de l'opération $\wedge$. \struck{\ill{}} C'est clair si on suppose qu'on est hors des pts de car.\ $2$, tous les facteurs étant alors étales, ou plus généralement, si tous sauf l'un sont étales --- mais c'est moins évident sans cette condition, même \uncertain{infinitésimale}. On aimerait une description directe d'un produit contracté \uncertain{multiple}
\[
  \bigwedge_{i \in I} X_i \qquad I \text{ un ens.\ fini}
\]
\struck{\ill{}} \ill{} en termes de \uncertain{formes} \uncertain{trilinéaires}. P.\ ex.\ sur $H_1$, $H_2$, $H_3$, formes \uncertain{trilinéaires} \struck{\ill{}} $f(u_1, u_2, u_3)$ telles que
\[
  \begin{aligned}
    f(1, u_2, u_3) &= \operatorname{Tr}(u_2)\operatorname{Tr}(u_3) \\
    f(u_1, 1, u_3) &= \operatorname{Tr}(u_1)\operatorname{Tr}(u_3) \\
    f(1, u_2, u_3) &= \operatorname{Tr}(u_2)\operatorname{Tr}(u_3)
  \end{aligned}
\]
\note{la troisième ligne répète la première telle qu'elle est écrite ; on attendrait $f(u_1, u_2, 1) = \operatorname{Tr}(u_1)\operatorname{Tr}(u_2)$.}
(elles forment \uncertain{bien} un torseur sous $(L_1 \otimes L_2 \otimes L_3)^{-1}$), satisfaisant à \ill{} conditions quadratiques convenables --- mais lesquelles ??

\section{Azumaya rg 4 (III)}
\note{titre de sa main en tête de la page 149, avec la date « Avril 86 » ; la page 148 est blanche. Un nouvel ensemble commence ici, numéroté III par lui.}

\page{149}
Avril 86

\textbf{\emph{Digression}} sur les produits tensoriels d'anneaux affines à involution.

La donnée, sur un $k$-module $H$ muni d'un élément $1_H$ (« unité ») \add{tel que $\lambda \mapsto \lambda 1$ soit injectif de $k$ dans $H$}, d'une \struck{involution} \add{appl.\ lin.}\ $x \mapsto x'$ qui a) fixe $1_H$ et b) induit sur $H/k1$ la symétrie, i.e.\ $x' \equiv -x \bmod k.1$ i.e.
\[
  x + x' \in k.1 ,
\]
équivaut à la donnée d'une forme linéaire $T : H \to k$ telles que $T(1) = 2$, et on aura
\[
  x' = T(x).1 - x \qquad \text{alors } x \mapsto x' \text{ est involution}
\]
\marginal{i.e.\ « une » $k \to H$}

Dualement, \struck{la donnée} \ill{} revient à se donner un $k$-module $E$ et une augmentation
\[
  E \xrightarrow{\ \varepsilon\ } k
\]
surjective (ce qui revient à se donner un espace affine $A = \varepsilon^{-1}(1)$ sur $k$, dont le module des translations sera $A^0 = \operatorname{Ker}(\varepsilon)$), \struck{revient à se donner une} la donnée d'une involution de $E$ compatible avec $\varepsilon$ (pour id.\ dans $k$) et induisant sur $A^0 = \operatorname{Ker}\varepsilon$ la symétrie $x \mapsto -x$, revient à se donner un élément
\[
  t \in E \quad \text{avec} \quad \varepsilon(t) = 2 ,
\]
l'involution sera alors
\[
  x \longmapsto \varepsilon(x)\,t - x
\]
et $t$ sera fixe. Se donner un espace affine $A$ sur $k$, muni d'une involution qui induise $-\mathrm{id}$ sur le module des translations, revient à se donner un tel espace\note{lecture de la fin de phrase incertaine.} (cet élément est alors invariant \ill{}).

\page{150}
Si $2$ est inv., il y a alors un seul pt fixe dans $A$, et $A$ devient ainsi un module, dit l'\uncertain{anneau} \ill{} que $x \mapsto -x$. Si par contre $2$ est \struck{inversible} \add{nul} \uncertain{dans} $k$, \struck{\ill{}} il \uncertain{s'ensuit} que $t$ est dans $k$…\note{lecture incertaine ; peut-être « $t$ est dans le module des translations ».} module des translations, et la donnée de $A$, et de $t \in A^0$ équivaut à celle de l'involution sur le translation par $t$. Nous \uncertain{nommerons} involutions affines « \uncertain{idéales} », \ill{} seules \uncertain{qui} \ill{} dans le cas où $k = \mathbb{Z}$, quand $2$ n'est ni inv., ni une div.\ \ill{} $k$. \struck{Je vais enfin}

Je voudrais définir une opération de produit tensoriel d'espaces affines à involutions. Soient $(A_1, A_2)$ deux tels espaces, avec involutions $\sigma_1$, $\sigma_2$ (je les notera\supplied{i} tous les deux $x \mapsto x'$, qui est ambigu \ill{}) et soit $B$ un espace affine sur $k$, les espaces de translations sont notés $A_1^0$, $A_2^0$, $B^0$.

\marginal{il \uncertain{faudrait} que $B$ soit un \struck{\ill{}} espace \struck{\ill{}} affine \ill{} \ill{}, ce qui \uncertain{conduirait} \ill{} \ill{} $B^0$ \uncertain{opérant}}
\note{note oblique dans la marge gauche, en partie biffée.}

\uncertain{Une} application
\[
  f : A_1 \times A_2 \longrightarrow B
\]
est dite $(\sigma_1, \sigma_2)$-bilinéaire, si \struck{\ill{}}
\[
  \text{\struck{$f(x$}} \quad \exists\, f^0 : A_1^0 \times A_2^0 \longrightarrow B^0 \quad \text{bil}
\]
\note{la page s'arrête sur cette ligne ; la condition se poursuit page 151.}

\page{151}
\note{suite de la condition commencée au bas de la page 150.}
tel qu'on ait
\[
  (*) \quad
  \left\{
  \begin{aligned}
    f(x_1 + \alpha_1, x_2) &= f(x_1, x_2) + f^0(\alpha_1, \sigma_2 x_2 - x_2) \\
    f(x_1, x_2 + \alpha_2) &= f(x_1, x_2) + f^0(\sigma_1 x_1 - x_1, \alpha_2)
  \end{aligned}
  \right.
\]
$\forall\, \alpha_1 \in A_1^0$, $\alpha_2 \in A_2^0$. On en conclut
\[
  \begin{aligned}
    f(x_1 + \alpha_1, x_2 + \alpha_2) = f(x_1, x_2) &+ f^0(\alpha_1, \sigma_2 x_2 - x_2) + f^0(\sigma_1 x_1 - x_1, \alpha_1) \\
    &- 2 f^0(\alpha_1, \alpha_2)
  \end{aligned}
\]
\note{dans le troisième terme on lit $\alpha_1$ en second argument ; on attendrait $\alpha_2$.}

Si on choisit $x_1$, $x_2$ comme origines, pour identifier $A_1$ à $A_1^0$ (via les $\alpha_1$) et $A_2$ à $A_2^0$ (via les $\alpha_2$), posant
\[
  \begin{aligned}
    \sigma_1 x_1 - x_1 &= t_1 - 2x_1 = \tau_1 \in A_1^0 \\
    \sigma_2 x_2 - x_2 &= t_2 - 2x_2 = \tau_2 \in A_2^0
  \end{aligned}
\]
la formule devient
\[
  \underset{\overset{\| \mathrm{df}}{f(x_1 + \alpha_1,\, x_2 + \alpha_2)}}{g(\alpha_1, \alpha_2)}
  = \underset{\overset{\|}{f(x_1, x_2)}}{c} + f^0(\alpha_1, \tau_2) + f^0(\tau_1, \alpha_2) - 2 f^0(\alpha_1, \alpha_2)
\]
\struck{\ill{}}. Cette formule définit entièrement la fonction $g$ (donc $f$) quand on connaît $c$ ($= f(x_1, x_2) = g(0,0)$), $\tau_1 \in A_1^0$, $\tau_2 \in A_2^0$ (qui définissent les involutions sur $A_1 \simeq A_1^0$, $A_2 \simeq A_2^0$ par $\alpha_1 \mapsto \tau_1 - \alpha_1$, $\alpha_2 \mapsto \tau_2 - \alpha_2$) via $x_1$, $x_2$, et l'application bil.\ $f^0$. Quand elles \ill{} \ill{} données, on \uncertain{vérifie} que la formule précédente définit bien une application $(\sigma_1, \sigma_2)$-bil. Mais si on ne suppose pas $2$ inv., il n'y a aucune

\marginal{Si $2.1_k = 0$, $g$ est \uncertain{vraiment} linéaire en $(\alpha_1, \alpha_2)$ et bien déterminée par $g(\alpha_1, \alpha_2) = c + f^0(\alpha_1, \tau_2) + f^0(\tau_1, \alpha_2)$. \quad Si $2$ inv., \ill{} \ill{} \ill{} \ill{} $\tau_1 = \tau_2 = 0$ \ill{} pour \ill{} $f^0(\alpha_1, \alpha_2)$ \ill{} $x_1'$, \ill{} \ill{} l'application bil.\ $f^0$ \ill{} \ill{}}
\note{deux notes obliques dans la marge gauche, en partie superposées.}

\page{152}
garantie que $f^0$ soit \ill{} \uncertain{par} $f$ l'est [Ex : si on prend $\tau' = \tau'' = 0$, $2.1_k = 0$, alors $f$ doit être \emph{constante}]. C'est pourquoi \uncertain{je} préfère inclure $f^0$ dans la structure de l'application --- i.e.\ l'application est une \uncertain{couple}
\[
  (f, f^0) : \quad
  \begin{aligned}
    f &: A_1 \times A_2 \longrightarrow B \\
    f^0 &: A_1^0 \times A_2^0 \longrightarrow B^0
  \end{aligned}
\]

On voit \struck{tout} de suite qu'il y a un accouplement \emph{universel}, dont le module des translations est
\[
  A_1^0 \otimes A_2^0
\]
et \uncertain{espace} noté\supplied{s} $A_1 * A_2$, espace affine sous $A_1^0 \otimes A_2^0$, avec $(\sigma_1, \sigma_2)$-accouplement
\[
  (x_1, x_2) \longmapsto x_1 * x_2 : A_1 \times A_2 \longrightarrow A_1 * A_2
\]
On peut le voir comme \ill{} sous-espace affine de
\[
  E_1 \otimes E_2 \supset A_1^0 \otimes A_2^0
\]
($E_i$ étant les \ill{} \uncertain{enveloppes} de l'espace affine $A_i$), l'application universelle étant induite par

\page{153}
\[
  \boxed{\;
  \begin{aligned}
    (x_1, x_2) \longmapsto x_1 * x_2 &= t_1 \otimes x_2 + x_1 \otimes t_2 - 2\, x_1 \otimes x_2 \\
    &= \sigma_1(x_1) \otimes x_2 + x_1 \otimes \sigma(x_2)
  \end{aligned}
  \;}
\]
\note{la seconde ligne est écrite telle quelle, sans indice à $\sigma(x_2)$.}
pendant que $f_{A_1, A_2}$, \uncertain{induite} sur $A_1^0 \times A_2^0$, est
\[
  \begin{aligned}
    f^0_{A_1, A_2} : A_1^0 \times A_2^0 &\longrightarrow A_1^0 \otimes A_2^0 \\
    (\alpha_1, \alpha_2) &\longmapsto \alpha_1 \otimes \alpha_2
  \end{aligned}
\]

\marginal{\ill{} de la \uncertain{fonction} \ill{} \ill{} \ill{} \uncertain{des} \ill{} $E_1 \otimes E_2$ par … $x_1 * x_2 = \sigma_1(x_1) \otimes x_2 + \ldots$ $= \varepsilon(x_1) t_1 \otimes x_2 + \varepsilon(x_2) \ldots$}
\note{note oblique dans l'angle supérieur gauche, très serrée.}

\struck{\ill{}} Bien sûr, $A_1 * A_2$ est un \uncertain{bifoncteur} en $A_1$, $A_2$. Donc il y a deux involutions, $\sigma_1 * \mathrm{id}_{A_2}$ et $\mathrm{id}_{A_1} * \sigma_2$ sur $A_1 * A_2$, définies par
\[
  x_1 * x_2 \longmapsto \sigma_1(x_1) * x_2 \qquad x_1 * x_2 \longmapsto x_1 * \sigma_2(x_2)
\]
Je dis qu'elles sont égales, et qu'on a
\[
  \underset{t_1 - x_1}{\underbrace{\sigma_1(x_1)}} * x_2 = x_1 * \underset{t_2 - x_2}{\underbrace{\sigma_2(x_2)}} = t_1 \otimes t_2 - x_1 * x_2
\]

\marginal{Si $2.1_k = 0$, $t_1 \otimes t_2 \in A_1^0 \otimes A_2^0$ \ill{} \ill{} translation \ill{} \ill{} $A_1 * A_2$ \ill{} \ill{} \ill{}}

Ainsi, $A_1 * A_2$ est \uncertain{muni} d'une involution $\sigma_{A_1 * A_2}$. Il est \uncertain{solution} du pb universel, cette fois dans les espaces affines involutifs sur $k$, pour les $\sigma$-accouplements d'espaces affines involutifs --- où cette fois on met la condition supplémentaire sur $f$
\[
  f(\sigma_1(x_1), x_2) = f(x_1, \sigma_2(x_2)) = \sigma_B f(x_1, x_2) .
\]

\page{154}
Ainsi, $A_1 * A_2$ est une opération interne dans les espaces affines à involutions. Je dis que cette opération \uncertain{admet} une structure (\uncertain{données} d'associativité ; \uncertain{commutativité}) \uncertain{avec} unité. La donnée est l'espace affine \uncertain{trivial} $k$ lui-même, avec \uncertain{comme} involution $x \mapsto 1 - x$ (\struck{donc \ill{}} c'est-à-dire \uncertain{l'analogue} \struck{\ill{}} de \uncertain{l'involution} triviale en car.\ $2$ !), l'enveloppe \struck{\ill{}} est \struck{\ill{} $\to k$}
\[
  k \times k \xrightarrow{\ \varepsilon\ } k , \qquad t_0 = (1,1) \mapsto 2
\]
($\sigma_0$ est $(\lambda, \mu) \mapsto (\mu, \lambda)$)
\[
  \begin{aligned}
    k &\hookrightarrow k \times k \\
    \lambda &\longmapsto (\lambda, 1 - \lambda)
  \end{aligned}
  \qquad
  \begin{aligned}
    k &\longleftarrow k \times k \\
    \lambda &\longmapsto (\lambda, -\lambda)
  \end{aligned}
\]
L'élément $e_0 = (1, 0)$ de $\operatorname{Ker} \varepsilon$\note{lecture du couple incertaine (« $(1,0)$ » surchargé).} \ill{} \uncertain{espace} $A_0$ satisfait
\[
  \sigma_0(e_0) - e_0 = 1 \in k = A_0^0
\]
Ceci dit, on utilise $e_0$ pour définir un \uncertain{isomorphisme} can.
\[
  \begin{aligned}
    A &\xrightarrow{\ \sim\ } A_0 * A \\
    x &\longmapsto e_0 * x
  \end{aligned}
\]
\note{après « $A_0 * A$ » une flèche suivie de quelques signes est biffée.}
compatible avec
\[
  A^0 \xrightarrow{\ \sim\ } \underset{\overset{\|}{A_0^0}}{k} \otimes A^0 \qquad \text{iso.\ can.}
\]

Pour définir la structure d'associativité, on va définir directement, pour un \uncertain{système} \uncertain{fini} d'espaces involutifs \uncertain{indexé} \uncertain{par} un \uncertain{ensemble} \uncertain{quelc.}\ de facteurs, le produit
\[
  A_1 * A_2 * \cdots * A_n
\]
\note{la page s'arrête ici ; la définition se poursuit page 155.}

\page{155}
\note{suite de la page 154.}
comme sol.\ d'un pb universel, que voici : on cherche, pour $A$ \uncertain{variant} \add{\uncertain{compact}}, \uncertain{un} espace affine $B$ sous un $B^0$, les applications $(f, f^0)$
\[
  \begin{aligned}
    f &: A_1 \times \cdots \times A_n \longrightarrow B \\
    f^0 &: A_1^0 \times \cdots \times A_n^0 \longrightarrow B^0
  \end{aligned}
\]
satisfaisant
\[
  \begin{aligned}
    f(x_1, \ldots, x_i + \alpha_i, \ldots, x_n) = f(x_1, \ldots, x_i, \ldots, x_n) \\
    + f^0(\sigma_1 x_1 - x_1, \ldots, \alpha_i, \ldots, \sigma_n x_n - x_n)
  \end{aligned}
\]
pour tout $i \in \{1, \ldots, n\}$ et $\alpha_i \in A_i^0$. Sans \uncertain{erreur}, la solution du pb universel se réalise comme sous-espace affine de
\[
  E_1 \otimes \cdots \otimes E_n
\]
\uncertain{sous} $E_1^0 \otimes \cdots \otimes E_n^0$\note{lecture des exposants incertaine : on attendrait $A_1^0 \otimes \cdots \otimes A_n^0$.}, par les éléments
\[
  x_1 * x_2 * \cdots * x_n \qquad (\text{pour } x_i \in A_i)
\]
définis ainsi
\[
  \begin{aligned}
    x_1 * x_2 * \cdots * x_n = \Bigl(\sum t_1 \cdots t_{n-1} x_n\Bigr) - 2\Bigl(\sum t_1 \cdots t_{n-2} x_{n-1} x_n\Bigr) \\
    + 4\Bigl(\sum t_1 \cdots t_{n-3} x_{n-2} x_{n-1} x_n\Bigr) + \cdots + (-1)^{n-1} 2^{n-1} x_1 x_2 \cdots x_n .
  \end{aligned}
\]
\note{le signe et le coefficient du dernier terme sont surchargés ; on lit $(-1)^{n-1}2^{n-1}$. Le signe entre les deux premières parenthèses est noirci ; on transcrit le $-$ que le calcul demande.}
Quand on change l'\emph{un} des $x_i$ en $\sigma_i(x_i) = t_i - x_i$, le produit est changé en
\[
  t_1 t_2 \cdots t_n - x_1 * x_2 * \cdots * x_n .
\]

\marginal{$*$ \ill{} \uncertain{engendré} « formule » \uncertain{(cf.\ d'abord} \ill{}) \quad $2.1_k = 0$, $t_i * x_i = 0$ \ill{} identique \ill{} \quad $x_1 * x_2 * \cdots * x_n = \ldots$ \quad Prob.\ \uncertain{que} \ill{} \ill{} \ill{} $\sigma_i(x_i)$ au \ill{} $x_i$}
\note{notes obliques dans la marge gauche, en partie illisibles.}

\page{156}
On peut quelques \uncertain{définir} \uncertain{multilinéairement} une application
\[
  \begin{aligned}
    E_1 \times \cdots \times E_n &\longrightarrow E_1 \otimes \cdots \otimes E_n \\
    (x_1, \ldots, x_n) &\longmapsto x_1 * x_2 * \cdots * x_n
  \end{aligned}
\]
définie par
\[
  \mathop{\ast}_i x_i = \sum_{i=1}^{n} (-1)^{i-1} 2^{i-1}
  \Bigl(\sum \underbrace{\varepsilon_1(x_1) t_1\, \varepsilon_2(x_2) t_2 \cdots \varepsilon_{n-i}(x_{n-i}) t_{n-i}}_{n-i \text{ facteurs } \varepsilon(x_j) t_j}\;
  \underbrace{x_{n-i+1} \cdots x_n}_{i \text{ facteurs}}\Bigr).
\]
Dont l'image est le sous-module \uncertain{engendré} par $\mathop{\ast} A_i$. Ce sous-module \uncertain{contient} \add{donc} $\bigotimes A_i^0$, et $t_1 t_2 \cdots t_n$, \uncertain{\ill{}} $\bigotimes A_i^0$ est \uncertain{engendré}, \add{$2$ !} \struck{ \ill{} engendré par $t_1 \cdots t_n$ (?)}. \marginal{\uncertain{voir} \ill{} \ill{}}

Formule \uncertain{plus} \uncertain{explicite} en termes des $\sigma_i$ :
\[
  \begin{aligned}
    x_1 * x_2 * \cdots * x_n &= \sum_{\substack{i \text{ impair} \\ 1 \leq i \leq n}} \ \sum_{\binom{n}{i} \text{ termes}} x_1 \cdots x_i\, \sigma_{i+1}(x_{i+1}) \cdots \sigma_n(x_n) && (2^{n-1} \text{ termes}) \\
    \sigma(x_1 * \cdots * x_n) &= \sum_{\substack{i \text{ pair} \\ 0 \leq i \leq n}} \ \sum_{\binom{n}{i} \text{ termes}} x_1 \cdots x_i\, \sigma_{i+1}(x_{i+1}) \cdots \sigma_n(x_n) && (2^{n-1} \text{ termes})
  \end{aligned}
\]
\[
  (x_1 * \cdots * x_n) + \sigma(x_1 * \cdots * x_n) = t_1 t_2 \cdots t_n\, \varepsilon_1(x_1) \varepsilon_2(x_2) \cdots \varepsilon_n(x_n)
\]
On a \struck{\ill{}}
\[
  (x_1 + \alpha_1) * x_2 * \cdots * x_n = \alpha_1 \otimes (\sigma_2(x_2) - x_2) \otimes \cdots \otimes (\sigma_n(x_n) - x_n)
\]
si $\alpha_1 \in E_1^0$, et \uncertain{formules} similaires.
\note{à gauche de « $(x_1 + \alpha_1)$ » un terme est noirci ; le membre de droite paraît omettre $x_1 * x_2 * \cdots * x_n +$, qui est peut-être sous la rature.}

Si $2$ est inversible, alors $t_1 \otimes \cdots \otimes t_n$ est le \uncertain{base} d'un supplémentaire canonique de $\bigotimes A_i^0$ dans le sous-module \add{de $\bigotimes E_i$} engendré par $\mathop{\ast}_i A_i$. On peut alors

\page{157}
considérer $E$ comme \uncertain{remplaçant} le module enveloppant de $\mathop{\ast} A_i$, en le \uncertain{munissant} de l'augmentation
\[
  \varepsilon : E \longrightarrow k
\]
définie par
\[
  \varepsilon(\underbrace{t_1 t_2 \cdots t_n}_{t}) = 2 .
\]
On notera que $\varepsilon$ n'\emph{est pas} induite par
\[
  \varepsilon_1 \otimes \cdots \otimes \varepsilon_n : \bigotimes E_i \longrightarrow k
\]
qui transforme $t_1 t_2 \cdots t_n$ non pas en $2$, mais en $2^n$.

Le groupe $(\mathbb{Z}/2\mathbb{Z})^n$ opère sur $P = \bigotimes E_i$ via les $\sigma_i$, et les différents caractères définissent des sous-modules propres, dont $\bigotimes E_i$ est la somme directe (si $2$ inv.\ !). Etant donné $\alpha = (\alpha_1, \ldots, \alpha_n) \in (\mathbb{Z}/2\mathbb{Z})^n$, soit
\[
  P_\alpha = \bigl\{ x \in P \bigm| \sigma_1^{\beta_1} \otimes \cdots \otimes \sigma_n^{\beta_n}(x) = (-1)^{\alpha_1\beta_1 + \alpha_2\beta_2 + \cdots + \alpha_n\beta_n} x \bigr\}
  \quad \forall \text{ syst.\ d'entiers } \beta_i \bmod 2 .
\]
Ceci dit, $E$ est le sous-espace propre associé au caractère $\alpha = (1, \ldots, 1)$ i.e.
\[
  E = P_{1, \ldots, 1} = \bigl\{ x \in P \bigm| (\sigma_1^{\beta_1} \otimes \cdots \otimes \sigma_n^{\beta_n})(x) = (-1)^{\sum \beta_i} x \bigr\}
\]
\note{la page s'arrête sur cette formule ; la suite est page 158.}

\page{158}
Finalement, on trouve des formules plus jolies en éliminant les $t_i$ et $\varepsilon_i$, et en se servant uniquement des $\sigma_i$. Cela nous amène à une \uncertain{variante} de \uncertain{l'objet} \uncertain{précédent} : celui des modules involutifs (i.e.\ \uncertain{modules} munis d'une involution). Eventuellement, on se donne dans $E_i$ les sous-modules $E_i^0$, de telle façon que $\sigma_i$ soit $-\mathrm{id}_{E_i^0}$ sur $E_i^0$, et soit id sur $E_i/E_i^0$ (car $E_i^0$ est unique si $2$ inv., donc c'est le sous-module propre $E_i^-$, formé des $x$ tels que $\sigma_i x = -x$). Un \emph{accouplement involutif} \struck{\ill{}} vers un module $E$, sera une application multilinéaire
\[
  \begin{aligned}
    E_1 \times \cdots \times E_n &\xrightarrow{\ f\ } E \\
    E_1^0 \times \cdots \times E_n^0 &\xrightarrow{\ f^0\ } E
  \end{aligned}
\]
\uncertain{munie} de \uncertain{satisfaisant} les relations de la forme
\[
  \begin{aligned}
    f(x_1 + \alpha_1, x_2, \ldots, x_n) = f(x_1, \ldots, x_n) + f^0(\alpha_1, \sigma_2(x_2) - x_2, \ldots \\
    \ldots, \sigma_n(x_n) - x_n)
  \end{aligned}
\]
pour $\alpha_1 \in E_1^0$ (et les $n - 1$ autres relations similaires). On \uncertain{aura} donc
\[
  f(\alpha_1, \alpha_2, \ldots, \alpha_n) = (-2)^{n-1} f^0(\alpha_1, \alpha_2, \ldots, \alpha_n)
\]
\note{la flèche $f^0$ est tracée de la seconde ligne vers le même $E$ ; « $f^0$ » est surchargé.}

\page{159}
La solution du pb universel est un \uncertain{module} \uncertain{contenant} \add{\uncertain{pas} \uncertain{exactement}} le \uncertain{module} $\bigotimes E_i^0 \oplus \bigotimes E_i$, \uncertain{extension} de $\bigotimes E_i/E_i^0$ par $\bigotimes E_i^0$ :
\[
  0 \longrightarrow \bigotimes E_i^0 \xrightarrow{\ i\ } \mathop{\ast}_i (E_i, \sigma_i)^{\wedge} \xrightarrow{\ J\ } \bigotimes (E_i/E_i^0) \longrightarrow 0
\]
\marginal{NB On \uncertain{dispose} \ill{} $E_i^-$ et $E_i/E_i^0$ \ill{} $E_i^+$ \ill{} $E_i$}

Quand on \uncertain{prend} \struck{\ill{}} \add{dans} $E_i$ \uncertain{tous} $E_i/E_i^0 \simeq k$, \struck{alors} \uncertain{ainsi} $\mathop{\ast}_i(E_i, \sigma_i)$ apparaît comme \uncertain{une} \ill{} par $\bigotimes E_i^0$ \ill{} \ill{} \ill{} le contexte des espaces affines. Nous \uncertain{trouvons} un hom.\ canonique
\[
  \mathop{\ast}_i (E_i, \sigma_i) \xrightarrow{\ \tau\ } \bigotimes E_i
\]
donné par la formule déjà écrite (pour $n = 2$, c'est la formule $x_1 * x_2 = \sigma_1 x_1 \otimes x_2 + x_1 \otimes \sigma_2 x_2$), mais \uncertain{on} \uncertain{a}

\begin{tikzcd}[column sep=large]
  \bigotimes E_i^0 \arrow[d, "i"'] \arrow[dr, "{(-2)^{n-1}\,\mathrm{can}}"] & \\
  \mathop{\ast}(E_i, \sigma_i) \arrow[r, "\tau"] \arrow[dr, "{2^{n-1}J}"'] & \bigotimes E_i \arrow[d, "\mathrm{can}"] \\
  & \bigotimes E_i/E_i^0
\end{tikzcd}

si $2$ \uncertain{inversible}, \uncertain{mais} \uncertain{pas} en général.

On a \uncertain{ainsi} une loi de composition (\uncertain{commutative}) \uncertain{avec} \uncertain{associativité}, \uncertain{commutation} \ill{}…

\marginal{\uncertain{Prendre} \ill{} \uncertain{pour} \ill{} « $E_i^0$ » \ill{} \uncertain{extension} \ill{} \ill{} $\bigotimes E_i/E_i^0$ \ill{}, \uncertain{donc} \struck{\ill{}} $\mathop{\ast}(E_i, \sigma_i) \simeq$ \ill{} $\sigma_i = \mathrm{id}$, \uncertain{retrouve} \ill{} \quad \uncertain{Peut-être} \ill{} \ill{} \ill{} 2 \ill{}}
\note{notes obliques dans la marge gauche, très serrées.}

\page{160}
Si $2$ inv., la catégorie des modules involutifs \uncertain{équivaut} à celle des \struck{\ill{}} couples de modules $(E^-, E^+)$ (avec $E^- = E^0$), et l'opération $*$ est l'opération de produit tensoriel \uncertain{signé} sur les deux facteurs --- aucun \uncertain{mystère} ! L'opération est délicate et utile, par contre, si $2$ non inversible (p.\ ex.\ en car.\ $2$).

Revenons au contexte des espaces affines. Si on a un espace affine $A$ sous un module de translations $L = A^0$ \emph{inversible}, alors les \uncertain{données} d'un diviseur quadratique dans $\check{\mathbb{V}}(A)$ \add{\uncertain{(ou} \uncertain{dans} $\check{\mathbb{V}}(A^0)$)} (la \uncertain{solution} \struck{\ill{}} définie par $A$) revient à la donnée d'une forme quadratique sur $E$ (\uncertain{enveloppe} affine de $A$)
\[
  q_E : E \longrightarrow L^{\otimes 2}
\]
\uncertain{i.e.}\ d'une fonction \uncertain{polynômiale} \add{\uncertain{sur} $A$} de degré $\leq 2$ sur $A$), telle que $q_L = q \mid L$ (la « partie homog.\ » \uncertain{de} $q$) soit l'\uncertain{élévation} au carré \uncertain{canonique}
\[
  L \longrightarrow L^{\otimes 2}, \qquad u \longmapsto u^{\otimes 2}
\]
\note{la page s'arrête ici ; l'argument se poursuit au-delà de ce lot.}

\end{document}
